% --- CORRECCIÓN IMPORTANTE EN LA PRIMERA LÍNEA ---
% Agregamos "xcolor=table" para que funcione \rowcolor sin errores
\documentclass[aspectratio=169, 10pt, xcolor=table]{beamer}

% --- Configuración del Tema (estilo profesional como el ejemplo) ---
\usetheme{Madrid}
\usecolortheme{dolphin}

% --- Paquetes necesarios ---
\usepackage[utf8]{inputenc}
\usepackage[spanish]{babel}
\usepackage{amsmath, amssymb}
\usepackage{booktabs}
\usepackage{graphicx}
\usepackage{listings}
\usepackage{tikz}
\usetikzlibrary{arrows.meta, positioning, shapes.geometric}

% --- Ruta de Imágenes (crea la carpeta "imagenes" y coloca ahí tus archivos) ---
\graphicspath{{./imagenes/}}

% --- Estilo para código Python ---
\definecolor{codegreen}{rgb}{0,0.6,0}
\definecolor{codegray}{rgb}{0.5,0.5,0.5}
\definecolor{codepurple}{rgb}{0.58,0,0.82}
\definecolor{backcolour}{rgb}{0.95,0.95,0.92}
\lstset{
    language=Python,
    backgroundcolor=\color{backcolour},
    commentstyle=\color{codegreen},
    keywordstyle=\color{blue},
    numberstyle=\tiny\color{codegray},
    stringstyle=\color{codepurple},
    basicstyle=\ttfamily\scriptsize,
    breaklines=true,
    frame=single,
    showstringspaces=false
}

% --- Pie de página personalizado ---
\makeatletter

\setbeamertemplate{footline}{
  \leavevmode%
  \hbox{%
  \begin{beamercolorbox}[wd=.333333\paperwidth,ht=2.25ex,dp=1ex,center]{author in head/foot}%
    \usebeamerfont{author in head/foot}\insertshortauthor
  \end{beamercolorbox}%
  \begin{beamercolorbox}[wd=.333333\paperwidth,ht=2.25ex,dp=1ex,center]{title in head/foot}%
    \usebeamerfont{title in head/foot}Sesión 1
  \end{beamercolorbox}%
  \begin{beamercolorbox}[wd=.333333\paperwidth,ht=2.25ex,dp=1ex,right]{date in head/foot}%
    \usebeamerfont{date in head/foot}NLP \hfill \insertframenumber/\inserttotalframenumber \hspace*{0.3cm}
  \end{beamercolorbox}}%
  \vskip0pt%
}

\makeatother

% --- Datos de la portada ---
\title[SESIÓN 1 NLP]{\textbf{Sesión 1}}
\subtitle{Introducción a NLP y Fundamentos Lingüísticos}
\author[Prof. C. Sánchez]{Carlos César Sánchez Coronel}
\institute{\textbf{NLP}}
\date{}

% Logo opcional (descomentar si tienes la imagen)
% \titlegraphic{\includegraphics[height=1.5cm]{logo.png}}

% --- Eliminar la barra de navegación (opcional) ---
\setbeamertemplate{navigation symbols}{}

% --- Índice al inicio de cada sección ---
\AtBeginSection[]{
  \begin{frame}
    \frametitle{Contenido}
    \tableofcontents[currentsection]
  \end{frame}
}

% ============================================================================
% INICIO DEL DOCUMENTO
% ============================================================================
\begin{document}

% --- Portada ---
\begin{frame}
    \titlepage
\end{frame}

% --- Índice general ---
\begin{frame}{Contenido}
    \tableofcontents
\end{frame}

% ============================================================================
% SECCIÓN 1: INTRODUCCIÓN A NLP
% ============================================================================
\section{Introducción a NLP}

\begin{frame}{¿Qué es NLP?}
    \begin{block}{Definición}
        Procesamiento del Lenguaje Natural (NLP) es una subdisciplina de la IA que se ocupa de dotar a las computadoras de la capacidad de entender, interpretar y generar lenguaje humano.
    \end{block}
    \begin{itemize}
        \item \textbf{Objetivo}: puente entre el lenguaje humano (discreto, simbólico) y las computadoras (números).
        \item \textbf{Desafíos únicos}: ambigüedad, creatividad, dependencia del contexto, conocimiento de mundo.
    \end{itemize}
    % Basado en introducción del PDF
\end{frame}

\begin{frame}{Breve historia del NLP}
    \begin{itemize}
        \item \textbf{1950-60}: enfoques basados en reglas, traducción palabra por palabra.
        \item \textbf{1980-90}: métodos estadísticos y corpus anotados.
        \item \textbf{2000}: modelos clásicos (Naive Bayes, SVM) aplicados a texto.
        \item \textbf{2010}: word embeddings (Word2Vec, GloVe) y redes neuronales.
        \item \textbf{2017-presente}: Transformers (BERT, GPT), modelos fundacionales.
    \end{itemize}
    % Basado en PDF
\end{frame}

% ============================================================================
% SECCIÓN 2: NIVELES DEL LENGUAJE
% ============================================================================
\section{Niveles del Lenguaje}

\begin{frame}{Morfología: estructura de las palabras}
    \begin{block}{Definición}
        Estudia la estructura interna de las palabras y sus constituyentes mínimos con significado: los \textbf{morfemas}.
    \end{block}
    \begin{exampleblock}{Ejemplo: ``inextinguible''}
        \begin{itemize}
            \item Prefijo ``in-'' (negación)
            \item Lexema ``extingu'' (raíz)
            \item Sufijo ``-ible'' (posibilidad)
        \end{itemize}
    \end{exampleblock}
    \textbf{Problemas computacionales:} segmentación morfológica, lematización, stemming.
    % Basado en PDF
\end{frame}

\begin{frame}{Sintaxis: estructura de las oraciones}
    \begin{block}{Definición}
        Estudia cómo se combinan las palabras para formar oraciones gramaticales.
    \end{block}
    \begin{itemize}
        \item \textbf{Estructura de constituyentes}: oración → SN + SV, SN → Det + N + Adj, etc.
        \item \textbf{Representación}: árboles sintácticos (grafo dirigido acíclico).
    \end{itemize}
    % Basado en PDF
    % IMAGEN: Árbol sintáctico simple
\end{frame}

\begin{frame}{Semántica: el significado}
    \begin{block}{Niveles semánticos}
        \begin{itemize}
            \item \textbf{Léxica}: significado de palabras y relaciones (sinonimia, antonimia, hiperonimia).
            \item \textbf{Composicional}: cómo se combinan los significados de las palabras.
            \item \textbf{Del discurso}: significado en contexto extendido.
        \end{itemize}
    \end{block}
    \begin{exampleblock}{Problemas computacionales}
        \begin{itemize}
            \item Desambiguación del sentido de las palabras (WSD)
            \item Reconocimiento de entidades nombradas (NER)
            \item Análisis de relaciones semánticas
        \end{itemize}
    \end{exampleblock}
    % Basado en PDF
\end{frame}

\begin{frame}{Pragmática: el uso en contexto}
    \begin{block}{Definición}
        Estudia cómo el contexto influye en la interpretación del significado: intención del hablante, conocimiento compartido, normas sociales.
    \end{block}
    \begin{exampleblock}{Ejemplos}
        \begin{itemize}
            \item Ironía: ``Qué buen tiempo hace'' durante una tormenta.
            \item Implicaturas: ``¿Tienes hora?'' no es una pregunta sobre posesión.
            \item Actos de habla: prometer, preguntar, ordenar.
        \end{itemize}
    \end{exampleblock}
    \textbf{Desafíos computacionales}: detección de sarcasmo, comprensión de intenciones.
    % Basado en PDF y Pregunta 1 (ironía)
\end{frame}

% ============================================================================
% SECCIÓN 3: PART-OF-SPEECH TAGGING
% ============================================================================
\section{Part-of-Speech (POS) Tagging}

\begin{frame}{Definición y formalización}
    \begin{block}{POS Tagging}
        Asignar a cada palabra de un texto su categoría gramatical (sustantivo, verbo, adjetivo, etc.) según su contexto.
    \end{block}
    \begin{block}{Formalización matemática}
        Dada una secuencia de palabras $w_1^n$, buscar la secuencia de etiquetas $t_1^n$ más probable:
        \[
        \hat{t}_1^n = \arg\max_{t_1^n} P(t_1^n | w_1^n) \propto P(w_1^n | t_1^n) P(t_1^n)
        \]
        Modelos: HMM (asume independencia de primer orden), CRF, BiLSTM, Transformers.
    \end{block}
    % Basado en PDF
\end{frame}

\begin{frame}{Conjuntos de etiquetas y aplicaciones}
    \begin{columns}[t]
        \column{0.5\textwidth}
        \textbf{Tagsets comunes}
        \begin{itemize}
            \item Penn Treebank (inglés): 45 etiquetas (NN, NNS, VB, VBD, JJ, RB...)
            \item EAGLES (lenguas europeas)
        \end{itemize}
        \column{0.5\textwidth}
        \textbf{Aplicaciones}
        \begin{itemize}
            \item Parsing sintáctico
            \item NER
            \item Análisis de sentimiento
        \end{itemize}
    \end{columns}
    \vspace{0.3cm}
    \begin{exampleblock}{Corpus anotado: Petrolés (Petro NLP)}
        Documentos técnicos geocientíficos en portugués con anotación POS para el dominio del petróleo y gas.
    \end{exampleblock}
    % Basado en PDF y Pregunta 2 (dominio médico)
\end{frame}

% ============================================================================
% SECCIÓN 4: PARSING (ANÁLISIS SINTÁCTICO)
% ============================================================================
\section{Parsing: Análisis Sintáctico}

\begin{frame}{Tipos de parsing}
    \begin{columns}[t]
        \column{0.5\textwidth}
        \textbf{Constituyentes}
        \begin{itemize}
            \item Árbol con frases (SN, SV)
            \item Gramáticas libres de contexto (CFG)
            \item Parsing probabilístico (PCFG)
        \end{itemize}
        \column{0.5\textwidth}
        \textbf{Dependencias}
        \begin{itemize}
            \item Relaciones binarias entre palabras
            \item Cada palabra tiene un ``cabeza'' y una relación
            \item Más útil para extracción de relaciones
        \end{itemize}
    \end{columns}
    \vspace{0.3cm}
    \begin{exampleblock}{Ejemplo de dependencias}
        ``Juan compró un auto'': \\
        \hspace*{1cm} ``compró'' es raíz, ``Juan'' es nsubj, ``auto'' es obj, ``un'' es det.
    \end{exampleblock}
    % Basado en PDF y Pregunta 4
\end{frame}

\begin{frame}{Formalización matemática (dependencias)}
    \begin{block}{Árbol de dependencias}
        Para una oración de $n$ palabras, $G=(V,E)$ con:
        \begin{itemize}
            \item $V = \{1,\dots,n\}$ (posiciones de las palabras)
            \item Cada vértice tiene exactamente un arco entrante, excepto la raíz.
            \item Grafo acíclico.
        \end{itemize}
    \end{block}
    \begin{block}{Objetivo}
        Maximizar la puntuación:
        \[
        T^* = \arg\max_T \sum_{(i,j)\in T} s(i,j)
        \]
        donde $s(i,j)$ es la puntuación de que $j$ dependa de $i$.
    \end{block}
    % Basado en PDF
\end{frame}

\begin{frame}{Treebanks (corpus de árboles sintácticos)}
    \begin{itemize}
        \item \textbf{CINTIL-TreeBank} (portugués): 10,039 oraciones, anotación semi-automática.
        \item \textbf{Penn Treebank} (inglés): estándar para constituyentes.
        \item \textbf{Universal Dependencies}: árboles de dependencias para múltiples idiomas.
    \end{itemize}
    % Basado en PDF
\end{frame}

% ============================================================================
% SECCIÓN 5: CORPUS Y DATASETS ANOTADOS
% ============================================================================
\section{Corpus y Datasets Anotados}

\begin{frame}{Tipos de corpus}
    \begin{itemize}
        \item \textbf{Generales}: British National Corpus.
        \item \textbf{Dominio específico}: medicina, derecho, petróleo.
        \item \textbf{Monolingües vs. multilingües}.
        \item \textbf{Anotados vs. no anotados}.
    \end{itemize}
    \vspace{0.3cm}
    \begin{block}{Principios FAIR}
        Findable, Accessible, Interoperable, Reusable. Esencial para la ciencia abierta.
    \end{block}
    % Basado en PDF y Pregunta 18
\end{frame}

\begin{frame}{Ejemplos de corpus anotados (I)}
    \begin{itemize}
        \item \textbf{RareDis} (salud): enfermedades raras y manifestaciones clínicas. 1,041 textos, F1 83.5\%.
        \item \textbf{SAMSum} (diálogos): 16,000 chats con resúmenes humanos.
        \item \textbf{ELEXIS} (multilingüe): 24 corpus (1,000M palabras cada uno) con anotación semántica.
    \end{itemize}
    % Basado en PDF
\end{frame}

\begin{frame}{Ejemplos de corpus anotados (II)}
    \begin{itemize}
        \item \textbf{Petro NLP} (industria petrolera): Petro KGraph, Petrolés, PetroGold, PetroNER, PetroRE.
        \item \textbf{Text2Story Lusa} (narrativas periodísticas): 357 artículos con anotación narrativa.
        \item \textbf{NLUCat} (asistentes virtuales): 12,000 instrucciones en catalán con intenciones y segmentos.
    \end{itemize}
    % Basado en PDF
\end{frame}

% ============================================================================
% SECCIÓN 6: APLICACIONES EN LA INDUSTRIA
% ============================================================================
\section{Aplicaciones en la Industria}

\begin{frame}{Chatbots y asistentes virtuales}
    \begin{itemize}
        \item Clasificación de intenciones, extracción de entidades, gestión de diálogo.
        \item Ejemplo: asistente bancario, NLUCat para personas vulnerables.
    \end{itemize}
    % Basado en PDF y Pregunta 20
\end{frame}

\begin{frame}{Análisis de sentimiento y redes sociales}
    \begin{itemize}
        \item Monitoreo de marca, detección de crisis.
        \item Desafío: ironía y sarcasmo (pragmática).
        \item Corpus con anotación de ironía.
    \end{itemize}
    % Basado en Pregunta 1
\end{frame}

\begin{frame}{Traducción automática y búsqueda semántica}
    \begin{itemize}
        \item Modelos secuencia a secuencia, Transformers.
        \item Búsqueda semántica con embeddings de oraciones (Sentence-BERT).
        \item Indexación con bases de datos vectoriales (FAISS, Pinecone).
    \end{itemize}
    % Basado en Preguntas 9 y 10
\end{frame}

\begin{frame}{Extracción de información en dominios específicos}
    \begin{itemize}
        \item Petro NLP: extracción de entidades geocientíficas de informes técnicos.
        \item RareDis: extracción de enfermedades y síntomas de literatura médica.
        \item Contratos legales: extracción de partes, fechas, cláusulas.
    \end{itemize}
    % Basado en PDF y Preguntas 6, 12
\end{frame}

% ============================================================================
% SECCIÓN 7: DESAFÍOS ACTUALES EN NLP
% ============================================================================
\section{Desafíos Actuales}

\begin{frame}{Ambigüedad y dependencia del contexto}
    \begin{itemize}
        \item \textbf{Ambigüedad}: ``banco'' (financiero / asiento).
        \item \textbf{Solución}: modelos contextuales como BERT.
    \end{itemize}
    % Basado en Pregunta 3
\end{frame}

\begin{frame}{Conocimiento de mundo y lenguas minoritarias}
    \begin{itemize}
        \item \textbf{Conocimiento de mundo}: ``El Papa no puede casarse'' requiere hechos extralingüísticos.
        \item \textbf{Solución}: integración con bases de conocimiento (RAG, Wikidata).
        \item \textbf{Lenguas minoritarias}: escasez de recursos. Transfer learning desde lenguas relacionadas.
    \end{itemize}
    % Basado en Preguntas 5 y 19
\end{frame}

\begin{frame}{Sesgos y equidad}
    \begin{itemize}
        \item Modelos aprenden sesgos de los datos.
        \item \textbf{Identificación}: comparar métricas por grupo demográfico.
        \item \textbf{Mitigación}: balanceo de datos, regularización de equidad, post-procesamiento.
        \item \textbf{Cumplimiento}: GDPR, Ley de IA de la UE.
    \end{itemize}
    % Basado en Pregunta 7
\end{frame}

% ============================================================================
% SECCIÓN 8: PREPARACIÓN PARA ENTREVISTAS
% ============================================================================
\section{Preparación para Entrevistas}

\begin{frame}{Preguntas conceptuales típicas}
    \begin{itemize}
        \item ¿Qué es NLP y cuáles son sus principales desafíos?
        \item Explica la diferencia entre sintaxis y semántica con ejemplos.
        \item ¿Qué es POS tagging y por qué es importante?
        \item ¿Qué es un corpus anotado? Menciona ejemplos.
    \end{itemize}
    % Basado en PDF
\end{frame}

\begin{frame}{Preguntas sobre aplicaciones y desambiguación}
    \begin{itemize}
        \item ¿Cómo diseñarías un chatbot para atención al cliente?
        \item ¿Cómo detectarías menciones de enfermedades en textos médicos?
        \item ¿Cómo resolverías la ambigüedad de ``banco''?
        \item ¿Qué consideraciones hay al crear un corpus de dominio específico?
    \end{itemize}
    % Basado en PDF y solucionario
\end{frame}

% ============================================================================
% SECCIÓN 9: RESUMEN
% ============================================================================
\section{Resumen}

\begin{frame}{Lo que hemos aprendido}
    \begin{exampleblock}{Conceptos clave}
        \begin{itemize}
            \item NLP: definición, historia, desafíos.
            \item Niveles del lenguaje: morfología, sintaxis, semántica, pragmática.
            \item POS tagging: formalización, modelos, aplicaciones, corpus.
            \item Parsing: constituyentes vs dependencias, treebanks.
            \item Corpus anotados: tipos, ejemplos (RareDis, SAMSum, Petro NLP, etc.), principios FAIR.
            \item Aplicaciones industriales: chatbots, sentimiento, traducción, extracción.
            \item Desafíos actuales: ambigüedad, conocimiento de mundo, lenguas minoritarias, sesgos.
            \item Preguntas típicas de entrevista.
        \end{itemize}
    \end{exampleblock}
\end{frame}

% --- Diapositiva final ---
\begin{frame}
    \centering
    \Huge \textbf{¡Gracias!}

\end{frame}

\end{document}